\documentclass[11pt,aspectratio=169]{beamer}
\usepackage[italian]{babel}
\usepackage[latin1]{inputenc}
\usepackage{graphicx}
\usepackage{listings}
\usepackage[export]{adjustbox}
% Custom bullets
\usepackage{pifont}
\usepackage[useregional]{datetime2}
\usepackage{tikz}
\usepackage[table]{colortbl}
\usepackage{venndiagram}
\usepackage{amsmath}

% custom legend
\newcommand{\cbox}[1]{\raisebox{\depth}{\fcolorbox{black}{#1}{\null}}}

\usetheme{Madrid}
\usecolortheme{spruce}

% Redefine labels
\deftranslation[to=Italian]{Section}{Sezione}
\deftranslation[to=Italian]{Subsection}{Sottosezione}

% Join commands
\def\ojoin{\setbox0=\hbox{$\bowtie$}%
  \rule[-.02ex]{.25em}{.4pt}\llap{\rule[\ht0]{.25em}{.4pt}}}
\def\leftouterjoin{\mathbin{\ojoin\mkern-5.8mu\bowtie}}
\def\rightouterjoin{\mathbin{\bowtie\mkern-5.8mu\ojoin}}
\def\fullouterjoin{\mathbin{\ojoin\mkern-5.8mu\bowtie\mkern-5.8mu\ojoin}}

\author[Piero Hierro \and Piermichele Rosati]{Piero Hierro \and Piermichele Rosati}

\institute[]{\large Universit\`a di Camerino\\ \footnotesize Tutorato - Basi di Dati}

\title[Esercizi - Algebra Relazionale e Fondamenti SQL]{3. Esercizi - Algebra Relazionale e Fondamenti SQL}
\setbeamertemplate{navigation symbols}{}
\setbeamertemplate{section in toc}[sections numbered]
\setbeamertemplate{subsection in toc}[subsections numbered]
%\titlegraphic{\includegraphics[width=6cm]{img/unicam-logo.jpg}}
\makeatletter
\setbeamertemplate{footline}{
    \leavevmode%
    \hbox{%
        \begin{beamercolorbox}[wd=.333333\paperwidth,ht=2.25ex,dp=1ex,center]{author in head/foot}%
            \usebeamerfont{author in head/foot}\insertshortauthor
        \end{beamercolorbox}%
        \begin{beamercolorbox}[wd=.333333\paperwidth,ht=2.25ex,dp=1ex,center]{title in head/foot}%
            \usebeamerfont{title in head/foot}\insertshorttitle
        \end{beamercolorbox}%
        \begin{beamercolorbox}[wd=.333333\paperwidth,ht=2.25ex,dp=1ex,right]{date in head/foot}%
            \usebeamerfont{date in head/foot}\insertshortdate\hspace*{5 em}
            \insertframenumber{} / \inserttotalframenumber\hspace*{2ex} 
        \end{beamercolorbox}%
    }%
    \vskip0pt%
}
\makeatother


\AtBeginSection[]
{
\begin{frame}<beamer>
\frametitle{Indice}
\tableofcontents[currentsection]
\end{frame}
}

\begin{document}
\begin{frame}
\centering
\includegraphics[width=5.5cm]{../img/unicam-logo.jpg}
\date{\today}
\titlepage
\end{frame}
\addtobeamertemplate{frametitle}{}{%
\begin{tikzpicture}[remember picture,overlay]
\node[anchor=north east,yshift=2pt] at (current page.north east) {\includegraphics[height=2cm]{../img/unicam-logo-notext.png}};
\end{tikzpicture}}

%\begin{frame}
%\titlepage
%\end{frame}

\section{Algebra Relazionale}
\def\schemaEx3.8{\small \begin{align*}
    DEPUTATI &= (\underline{Codice}, Cognome, Nome, Commissione, Provincia, Collegio)\\
    COLLEGI &= (\underline{Provincia}, Numero,Nome)\\
    PROVINCE &= (\underline{Sigla}, Nome, Regione)\\
    REGIONI &= (\underline{Codice}, Nome)\\
    COMMISSIONI &= (\underline{Numero}, Nome, Presidente)\\
    \end{align*}}
\begin{frame}{Esercizio 3.8}
    \framesubtitle{3.8.1}
    \schemaEx3.8
    Trovare nome e cognome dei presidenti di commissioni cui partecipa almeno un deputato eletto in una provincia della Sicilia.
\end{frame}
%
\begin{frame}{Esercizio 3.8}
    \framesubtitle{Soluzione 3.8.1}
    \vspace*{-1.2cm}
    \schemaEx3.8
    {\small Nome e cognome dei presidenti di commissioni cui partecipa almeno un deputato eletto in una provincia della Sicilia:}
    \small
    \begin{gather*}
        \onslide<8->{\pi_{Nom,Cogn}}\\
        \onslide<7->{((\rho_{Nom,Cogn~\leftarrow~Nome,Cognome}}
        \onslide<6->{(DEPUTATI))~\bowtie_{Presidente~=~Codice}}\\
        \onslide<5->{(COMMISSIONI~\bowtie_{Numero~=~Comm}}
        \onslide<4->{(\rho_{Comm~\leftarrow~{Commissione}}}
        \onslide<3->{(DEPUTATI~\bowtie_{Provincia~=~Sigla}}\\
        \onslide<2->{(PROVINCE~\bowtie_{Regione~=~Codice}}\\
        \onslide<1->{(\sigma_{Nome~=~'Sicilia'}(REGIONI))}
        \onslide<2->{)}
        \onslide<3->{)}
        \onslide<4->{)}
        \onslide<5->{)}
        \onslide<6->{)}
    \end{gather*}
\end{frame}
%
\begin{frame}{Esercizio 3.8}
    \framesubtitle{3.8.2}
    \schemaEx3.8
    Trovare nome e cognome dei deputati della commissione Bilancio.
\end{frame}
%
\begin{frame}{Esercizio 3.8}
    \framesubtitle{Soluzione 3.8.2}
    \vspace*{-1.2cm}
    \schemaEx3.8
    {\small Nome e cognome dei deputati della commissione Bilancio:}
    \small
    \begin{gather*}
        \onslide<4->{\pi_{NomeC,Cognome}(}
        \onslide<3->{(\rho_{NomeC~\leftarrow~Nome}}
        \onslide<2->{(DEPUTATI))~\bowtie_{Commissione~=~Numero}}\\
        \onslide<1->{(\sigma_{Nome~=~'Bilancio'}(COMMISSIONI))}
    \end{gather*}
\end{frame}
%
\begin{frame}{Esercizio 3.8}
    \framesubtitle{3.8.3}
    \schemaEx3.8
    Trovare nome, cognome e provincia di elezione dei deputati della commissione Bilancio.
\end{frame}
%
\begin{frame}{Esercizio 3.8}
    \framesubtitle{Soluzione 3.8.3}
    \vspace*{-1.2cm}
    \schemaEx3.8
    {\small Nome, cognome e provincia di elezione dei deputati della commissione Bilancio:}
    \small
    \begin{gather*}
        \onslide<6->{\pi_{NomeC,Cognome,NomeP}(}\\
        \onslide<5->{(\rho_{NomeP~\leftarrow~Nome}}
        \onslide<4->{(PROVINCE))~\bowtie_{Sigla~=~Provincia}}\\
        \onslide<3->{((\rho_{NomeC~\leftarrow~Nome}}
        \onslide<2->{(DEPUTATI))~\bowtie_{Commissione~=~Numero}}\\
        \onslide<1->{(\sigma_{Nome~=~'Bilancio'}(COMMISSIONI))}
        \onslide<3->{)}
    \end{gather*}
\end{frame}
%
\begin{frame}{Esercizio 3.8}
    \framesubtitle{3.8.4}
    \schemaEx3.8
    Trovare nome, cognome, provincia e regione di elezione dei deputati della commissione Bilancio.
\end{frame}
%
\begin{frame}{Esercizio 3.8}
    \framesubtitle{Soluzione 3.8.4}
    \vspace*{-1.2cm}
    \schemaEx3.8
    {\small Nome, cognome, provincia e regione di elezione dei deputati della commissione Bilancio:}
    \small
    \begin{gather*}
        \onslide<8->{\pi_{NomeC,Cognome,NomeP,NomeR}(}\\
        \onslide<7->{(\rho_{NomeR~\leftarrow~Nome}}
        \onslide<6->{(REGIONI))~\bowtie_{Codice~=~Regione}}\\
        \onslide<5->{(\rho_{NomeP~\leftarrow~Nome}}
        \onslide<4->{(PROVINCE))~\bowtie_{Sigla~=~Provincia}}\\
        \onslide<3->{((\rho_{NomeC~\leftarrow~Nome}}
        \onslide<2->{(DEPUTATI))~\bowtie_{Commissione~=~Numero}}\\
        \onslide<1->{(\sigma_{Nome~=~'Bilancio'}(COMMISSIONI))}
        \onslide<3->{)}
    \end{gather*}
\end{frame}
%
\section{Fondamenti SQL}
%%%%%%%%%%%%%%%%
\begin{frame}{GROUP BY}
\begin{table}[h]
\centering
\begin{minipage}{.45\textwidth}
\centering
\begin{tabular}{|c|c|c|}
\hline
\rowcolor{cyan!30} \multicolumn{3}{|c|}{Esame} \\
\hline
\rowcolor{cyan!30} matricola  & voto & corso \\
\hline
100  & 30 & 27035 \\
100  & 28 & 27038 \\
100  & 18 & 27010 \\
102  & 21 & 27038 \\
101  & 25 & 27045 \\
102 & 21 & 27010 \\
\hline
\end{tabular}
\end{minipage}%
\begin{minipage}{.45\textwidth}
\centering
\begin{tabular}{|c|c|}
\hline
\rowcolor{cyan!30} \multicolumn{2}{|c|}{Result} \\
\hline
\rowcolor{cyan!30} matricola  & voti \\
\hline
100  & 3 \\
101  & 1  \\
102  & 2 \\
\hline
\end{tabular}
\end{minipage}
\end{table}
\vspace{2em}
\texttt{SELECT matricola, COUNT(voto) AS voti\\FROM Esame\\GROUP BY matricola;}
\end{frame}
%%%%%%%%%%%%%%%%
\begin{frame}{HAVING}
\begin{table}[h]
\centering
\begin{minipage}{.45\textwidth}
\centering
\begin{tabular}{|c|c|c|}
\hline
\rowcolor{cyan!30} \multicolumn{3}{|c|}{Esame} \\
\hline
\rowcolor{cyan!30} matricola  & voto & corso \\
\hline
100  & 30 & 27035 \\
100  & 28 & 27038 \\
100  & 18 & 27010 \\
102  & 21 & 27038 \\
101  & 25 & 27045 \\
102 & 21 & 27010 \\
\hline
\end{tabular}
\end{minipage}%
\begin{minipage}{.45\textwidth}
\centering
\begin{tabular}{|c|c|c|}
\hline

\hline
\end{tabular}
\end{minipage}
\end{table}
\vspace{2em}
\texttt{SELECT matricola, COUNT(voto) AS voti\\FROM Esame\\GROUP BY matricola\\HAVING COUNT(voto)>1;}
\end{frame}
%%%%%%%%%%%%%%%%
\begin{frame}{HAVING}
\begin{table}[h]
\centering
\begin{minipage}{.45\textwidth}
\centering
\begin{tabular}{|c|c|c|}
\hline
\rowcolor{cyan!30} \multicolumn{3}{|c|}{Esame} \\
\hline
\rowcolor{cyan!30} matricola  & voto & corso \\
\hline
100  & 30 & 27035 \\
100  & 28 & 27038 \\
100  & 18 & 27010 \\
102  & 21 & 27038 \\
101  & 25 & 27045 \\
102 & 21 & 27010 \\
\hline
\end{tabular}
\end{minipage}%
\begin{minipage}{.45\textwidth}
\centering
\begin{tabular}{|c|c|}
\hline
\rowcolor{cyan!30} \multicolumn{2}{|c|}{Result} \\
\hline
\rowcolor{cyan!30} matricola  & voti \\
\hline
100  & 3 \\
102  & 2  \\
\hline
\end{tabular}
\end{minipage}
\end{table}
\vspace{2em}
\texttt{SELECT matricola, COUNT(voto) AS voti\\FROM Esame\\GROUP BY matricola\\HAVING COUNT(voto)>1;}
\end{frame}
%%%%%%%%%%%%%%%%
\begin{frame}{ORDER BY}
\begin{table}[h]
\centering
\begin{minipage}{.45\textwidth}
\centering
\begin{tabular}{|c|c|c|}
\hline
\rowcolor{cyan!30} \multicolumn{3}{|c|}{Esame} \\
\hline
\rowcolor{cyan!30} matricola  & voto & corso \\
\hline
100  & 30 & 27035 \\
100  & 28 & 27038 \\
100  & 18 & 27010 \\
102  & 21 & 27038 \\
101  & 25 & 27045 \\
102 & 21 & 27010 \\
\hline
\end{tabular}
\end{minipage}%
\begin{minipage}{.45\textwidth}
\centering
\begin{tabular}{|c|c|c|}
\hline

\hline
\end{tabular}
\end{minipage}
\end{table}
\vspace{2em}
\texttt{SELECT *\\FROM Esame\\ORDER BY matricola ASC;}
\end{frame}
%%%%%%%%%%%%%%%%
\begin{frame}{ORDER BY}
\begin{table}[h]
\centering
\begin{minipage}{.45\textwidth}
\centering
\begin{tabular}{|c|c|c|}
\hline
\rowcolor{cyan!30} \multicolumn{3}{|c|}{Esame} \\
\hline
\rowcolor{cyan!30} matricola  & voto & corso \\
\hline
100  & 30 & 27035 \\
100  & 28 & 27038 \\
100  & 18 & 27010 \\
102  & 21 & 27038 \\
101  & 25 & 27045 \\
102 & 21 & 27010 \\
\hline
\end{tabular}
\end{minipage}%
\begin{minipage}{.45\textwidth}
\centering
\begin{tabular}{|c|c|c|}
\hline
\rowcolor{cyan!30} \multicolumn{3}{|c|}{Esame} \\
\hline
\rowcolor{cyan!30} matricola  & voto & corso \\
\hline
100 & 30 & 27035 \\
100 & 28 & 27038 \\
100 & 18 & 27010 \\
101 & 25 & 27045 \\
102 & 21 & 27038 \\
102 & 21 & 27010 \\
\hline
\end{tabular}
\end{minipage}
\end{table}
\vspace{2em}
\texttt{SELECT *\\FROM Esame\\ORDER BY matricola ASC;}
\end{frame}
%%%%%%%%%%%%%%%%
\begin{frame}{ORDER BY}
\begin{table}[h]
\centering
\begin{minipage}{.45\textwidth}
\centering
\begin{tabular}{|c|c|c|}
\hline
\rowcolor{cyan!30} \multicolumn{3}{|c|}{Esame} \\
\hline
\rowcolor{cyan!30} matricola  & voto & corso \\
\hline
100  & 30 & 27035 \\
100  & 28 & 27038 \\
100  & 18 & 27010 \\
102  & 21 & 27038 \\
101  & 25 & 27045 \\
102 & 21 & 27010 \\
\hline
\end{tabular}
\end{minipage}%
\begin{minipage}{.45\textwidth}
\centering
\begin{tabular}{|c|c|c|}
\hline

\hline
\end{tabular}
\end{minipage}
\end{table}
\vspace{2em}
\texttt{SELECT *\\FROM Esame\\ORDER BY matricola DESC, voto ASC;}
\end{frame}
%%%%%%%%%%%%%%%%
\begin{frame}{ORDER BY}
\begin{table}[h]
\centering
\begin{minipage}{.45\textwidth}
\centering
\begin{tabular}{|c|c|c|}
\hline
\rowcolor{cyan!30} \multicolumn{3}{|c|}{Esame} \\
\hline
\rowcolor{cyan!30} matricola  & voto & corso \\
\hline
100  & 30 & 27035 \\
100  & 28 & 27038 \\
100  & 18 & 27010 \\
102  & 21 & 27038 \\
101  & 25 & 27045 \\
102 & 21 & 27010 \\
\hline
\end{tabular}
\end{minipage}%
\begin{minipage}{.45\textwidth}
\centering
\begin{tabular}{|c|c|c|}
\hline
\rowcolor{cyan!30} \multicolumn{3}{|c|}{Esame} \\
\hline
\rowcolor{cyan!30} matricola  & voto & corso \\
\hline
102 & 21 & 27038 \\
102 & 21 & 27010 \\
101 & 25 & 27045 \\
100 & 18 & 27010 \\
100 &  28 & 27038 \\
100 & 30 & 27035 \\
\hline
\end{tabular}
\end{minipage}
\end{table}
\vspace{2em}
\texttt{SELECT *\\FROM Esame\\ORDER BY matricola DESC, voto ASC;}
\end{frame}
%%%%%%%%%%%%%%%%
\begin{frame}{UNION}
\begin{table}[h]
\centering
\begin{minipage}{.45\textwidth}
\centering
\begin{tabular}{|c|c|c|}
\hline
\rowcolor{cyan!30} \multicolumn{3}{|c|}{Studente} \\
\hline
\rowcolor{cyan!30} Nome  & Cognome & AnnoDiNascita \\
\hline
ciccio  & pasticcio & 1992 \\
pippo  & calippo & 1994 \\

\hline
\end{tabular}
\end{minipage}%
\begin{minipage}{.45\textwidth}
\centering
\begin{tabular}{|c|c|c|}
\hline
\rowcolor{cyan!30} \multicolumn{3}{|c|}{Professore} \\
\hline
\rowcolor{cyan!30} Nome  & Cognome & DataNascita \\
\hline
balema & giuseppina & 1967 \\

\hline
\end{tabular}
\end{minipage}
\end{table}
\vspace{2em}
\texttt{SELECT * FROM Studente\\UNION\\SELECT * FROM Professore;}
\end{frame}
%%%%%%%%%%%%%%%%
\begin{frame}{UNION}
\begin{table}[h]
\centering
\begin{minipage}{.45\textwidth}
\centering
\begin{tabular}{|c|c|c|}
\hline
\rowcolor{cyan!30} \multicolumn{3}{|c|}{Studente} \\
\hline
\rowcolor{cyan!30} Nome  & Cognome & AnnoDiNascita \\
\hline
ciccio  & pasticcio & 1992 \\
pippo  & calippo & 1994 \\

\hline
\end{tabular}
\end{minipage}%
\begin{minipage}{.45\textwidth}
\centering
\begin{tabular}{|c|c|c|}
\hline
\rowcolor{cyan!30} \multicolumn{3}{|c|}{Professore} \\
\hline
\rowcolor{cyan!30} Nome  & Cognome & AnnoDiNascita \\
\hline
balema & giuseppina & 1967 \\

\hline
\end{tabular}
\end{minipage}
\end{table}
\vspace{2em}
\texttt{SELECT * FROM Studente\\UNION\\SELECT * FROM Professore;}
\begin{itemize}
    \item Stesso numero di colonne nelle SELECT
    \item Stesso ordine di colonne nelle SELECT
    \item Stesso tipo dato nelle colonne delle SELECT
\end{itemize}
\end{frame}
%%%%%%%%%%%%%%%%
\end{document}
   
\begin{frame}{Esercizio 4.1}
         Ordinare i seguenti domini in base al valore massimo rappresentabile, supponendo che \texttt{integer} abbia una rappresentazione a 32 bit e \texttt{smallint} a 16 bit:
        \begin{enumerate}
            \item \texttt{numeric(12,4)}
            \item \texttt{decimal(10)}
            \item \texttt{decimal(9)}
            \item \texttt{integer}
            \item \texttt{smallint}
            \item \texttt{decimal(6,1)}
        \end{enumerate}
\end{frame}

\begin{frame}{Esercizio 4.1}
    \textbf{Soluzione:}
    \begin{table}[h]
        \centering
        \begin{tabular}{|c|c|}
            \hline
            Dominio & Valore Massimo \\
            \hline
            \texttt{decimal(10)} & 9999999999 \\
            \texttt{integer} & 4294967296 \\
            \texttt{decimal(9)} & 999999999 \\
            \texttt{numeric(12,4)} & 99999999.9999 \\
            \texttt{decimal(6,1)} & 99999.9 \\
            \texttt{smallint} & 65536 \\
            \hline
        \end{tabular}
    \end{table}
\end{frame}
%%%%%%%%%%%%%%%%
\begin{frame}{Esercizio 4.7}
    Con riferimento ad una relazione \texttt{PROFESSORI(CF, Nome, Et\`a, Qualifica)}, scrivere le interrogazioni SQL che calcolano l'et\`a media dei professori di ciascuna qualifica, nei due casi seguenti:
    \begin{enumerate}
        \item se l'et\`a non \`e nota si usa per essa il valore nullo
        \item se l'et\`a non \`e nota si usa per essa il valore 0.
    \end{enumerate}
\end{frame}

\begin{frame}{Esercizio 4.7}
    \textbf{Soluzione:}
    Le funzioni aggregative escludono dalla valutazione le ennuple con valori nulli:
    \vspace{1em}
    
    \texttt{SELECT Qualifica, AVG(Et\`a) AS Et\`aMedia\\FROM Professori\\GROUP BY Qualifica;}
\end{frame}

\begin{frame}{Esercizio 4.7}
    \textbf{Soluzione:}
    \`E necessario escludere esplicitamente dal calcolo della media le ennuple con il valore che denota l'informazione incompleta:
    \vspace{1em}
    
    \texttt{SELECT Qualifica, AVG(Et\`a) AS Et\`aMedia\\FROM Professori\\WHERE Et\`a <> 0\\GROUP BY Qualifica;}
\end{frame}
%%%%%%%%%%%%%%%%
\begin{frame}{Esercizio 4.14}
    Dato il seguente schema:
    \begin{itemize}
        \item AEROPORTO(Citt\`a, Nazione,NumPiste)
        \item VOLO(IdVolo,GiornoSett,Citt\`aPart,OraPart,
        Citt\`aArr,OraArr,TipoAereo)
        \item AEREO(TipoAereo,NumPasseggeri,QtaMerci)
    \end{itemize}
    \vspace{1em}
    
    \textbf{Le citt\`a con un aeroporto di cui non \`e noto il numero di piste.}
\end{frame}

\begin{frame}{Esercizio 4.14}
    \textbf{Soluzione:}
    \vspace{1em}
    
    \texttt{SELECT DISTINCT Citt\`a\\FROM AEROPORTO\\WHERE NumPiste IS NULL;}
    \end{frame}
%%%%%%%%%%%%%%%%%%%
\begin{frame}{Esercizio 4.14}
    Dato il seguente schema:
    \begin{itemize}
        \item AEROPORTO(Citt\`a, Nazione,NumPiste)
        \item VOLO(IdVolo,GiornoSett,Citt\`aPart,OraPart,
        Citt\`aArr,OraArr,TipoAereo)
        \item AEREO(TipoAereo,NumPasseggeri,QtaMerci)
    \end{itemize}
    \vspace{1em}
    
    \textbf{Le nazioni da cui parte e arriva il volo con codice AZ274.}
\end{frame}

\begin{frame}{Esercizio 4.14}
    \textbf{Soluzione:}
    \vspace{1em}
    
    \texttt{SELECT A1.Nazione, A2.Nazione\\FROM AEROPORTO AS A1\\ JOIN VOLO ON A1.Citt\`a = Citt\`aArr \\JOIN AEROPORTO AS A2 ON Citt\`aPart = A2.Citt\`a\\WHERE IdVolo = 'AZ274';
}
    \end{frame}
%%%%%%%%%%%%%%%%%%%
\begin{frame}{Esercizio 4.14}
    Dato il seguente schema:
    \begin{itemize}
        \item AEROPORTO(Citt\`a, Nazione,NumPiste)
        \item VOLO(IdVolo,GiornoSett,Citt\`aPart,OraPart,
        Citt\`aArr,OraArr,TipoAereo)
        \item AEREO(TipoAereo,NumPasseggeri,QtaMerci)
    \end{itemize}
    \vspace{1em}
    
    \textbf{Le citt\`a con un aeroporto di cui non \`e noto il numero di piste.}
\end{frame}

\begin{frame}{Esercizio 4.14}
    \textbf{Soluzione:}
    \vspace{1em}
    
    \texttt{SELECT DISTINCT Citt\`a\\FROM AEROPORTO\\WHERE NumPiste IS NULL;}
    \end{frame}
%%%%%%%%%%%%%%%%%%%
\begin{frame}{Esercizio 4.14}
    Dato il seguente schema:
    \begin{itemize}
        \item AEROPORTO(Citt\`a, Nazione,NumPiste)
        \item VOLO(IdVolo,GiornoSett,Citt\`aPart,OraPart,
        Citt\`aArr,OraArr,TipoAereo)
        \item AEREO(TipoAereo,NumPasseggeri,QtaMerci)
    \end{itemize}
    \vspace{1em}
    
    \textbf{I tipi di aereo usati nei voli che partono da Torino}
\end{frame}

\begin{frame}{Esercizio 4.14}
    \textbf{Soluzione:}
    \vspace{1em}
    
    \texttt{SELECT TipoAereo\\FROM VOLO\\WHERE Citt\`aPart='Torino'}
    \end{frame}
%%%%%%%%%%%%%%%%%%%
\begin{frame}{Esercizio 4.14}
    Dato il seguente schema:
    \begin{itemize}
        \item AEROPORTO(Citt\`a, Nazione,NumPiste)
        \item VOLO(IdVolo,GiornoSett,Citt\`aPart,OraPart,
        Citt\`aArr,OraArr,TipoAereo)
        \item AEREO(TipoAereo,NumPasseggeri,QtaMerci)
    \end{itemize}
    \vspace{1em}
    
    \textbf{Le citt\`a con un aeroporto di cui non \`e noto il numero di piste.}
\end{frame}

\begin{frame}{Esercizio 4.14}
    \textbf{Soluzione:}
    \vspace{1em}
    
    \texttt{SELECT DISTINCT Citt\`a\\FROM AEROPORTO\\WHERE NumPiste IS NULL;}
    \end{frame}
%%%%%%%%%%%%%%%%%%%
\begin{frame}{Esercizio 4.14}
    Dato il seguente schema:
    \begin{itemize}
        \item AEROPORTO(Citt\`a, Nazione,NumPiste)
        \item VOLO(IdVolo,GiornoSett,Citt\`aPart,OraPart,
        Citt\`aArr,OraArr,TipoAereo)
        \item AEREO(TipoAereo,NumPasseggeri,QtaMerci)
    \end{itemize}
    \vspace{1em}
    
    \textbf{I tipi di aereo e il corrispondente numero di passeggeri per i tipi di aereo usati nei voli che partono da Torino. Se la descrizione dell`aereo non \`e disponibile, visualizzare solamente il tipo;}
\end{frame}

\begin{frame}{Esercizio 4.14}
    \textbf{Soluzione:}
    \vspace{1em}
    
    \texttt{SELECT VOLO.TipoAereo, NumPasseggeri\\FROM VOLO LEFT JOIN AEREO ON VOLO.TipoAereo=aereo.TipoAereo\\WHERE Citt\`aPart= 'Torino'}
    \end{frame}
   
%
\section{Conclusioni}

\begin{frame}{Domande?}
    \begin{figure}
\centering
    \includegraphics[width=0.75\textwidth]{../img/questions.jpg}
\end{figure}
\end{frame}

\begin{frame}{Fine}
    \centering
    \huge Grazie dell'attenzione!
\end{frame}

\end{document}